\section{Conclusion}\label{sec:conclusion}

This study compared Kubernetes configurations generated by five current LLMs with the manifests shipped by the same seven open-source projects, deploying and exercising every configuration rather than only scanning it. Each model received only the project's own documentation, under a neutral and a production-oriented prompt, three times each.

First, generated configurations are largely valid: 92.9\% are schema-valid, and requirement satisfaction is high for every model. Under a neutral prompt their static security is close to what these projects ship (within about ten Kubescape points per model), and GPT's and Claude's cover more of the reliability checklist. Validity does not imply a working deployment, however: only 59.6\% of the generated configurations became fully Ready and passed their smoke test, and higher static security and reliability scores showed no positive association with success. Second, asking for production-ready output raises static scores while lowering deployment success: across model--project cells, median Kubescape compliance rose by 9 points and checklist coverage by 22.5, while the smoke-test pass rate fell from 76.8\% to 60.6\%. The added failures come from hardening applied without regard to the image (capability drops that stop databases from starting, liveness probes that kill slow-starting services). Third, models translate a project's Docker Compose file too literally, carrying over host-level settings such as a mounted Docker socket that the maintainers deliberately left out of their Kubernetes manifests.

A generated configuration should therefore be deployed and exercised before it is trusted, and its hardening settings checked against the images they constrain. Future work should extend the evaluation to more projects, to managed multi-node clusters and to faults that exercise configuration choices (node drains that test disruption budgets, resource pressure, dependency restarts), and should test whether the cluster's record of a failed deployment, which identified a cause for nearly every failure here, lets a model repair its own configuration.
